\documentclass{exam}
\usepackage{../../shah350}

\hypersetup{colorlinks=true, linktoc=section, linkcolor=blue}

\pagestyle{headandfoot}
\firstpageheadrule
\runningheadrule
\firstpageheader{Prof. Lepowsky \\ Linear Algebra}{Homework \#8}{\href{https://github.com/jeevanshah07/MATH350/blob/main/homework/homework8/main.tex}{\underline{Source}} \\ Jeevan Shah}
\runningheader{Linear Algebra \\ Homework \#8}{}{Shah}
\firstpagefooter{}{\thepage}{}
\runningfooter{ }{\thepage}{ }

\printanswers

% 5.1.4c,d; 5.1.5f; 5.1.9; 5.1.10; 5.2.2d,f; 5.2.3d; 5.2.13; 5.2.21; 5.2.22

\begin{document}
  \begin{problem}{5.1.4}
    \begin{parts}
      \setcounter{partno}{2}
      \part \begin{enumerate}[label=\roman*.)]
        \item The characteristic polynomial for $A$ is 
        \[
          \lambda^2 - 1 = 0 \Rightarrow \lambda = \pm 1
        \]
        \item For $\lambda = 1$: 
        \[
          A - I = \begin{pmatrix}
            i - 1 & 1 \\
            2 & -1 - i
          \end{pmatrix}
          \Rightarrow \begin{cases}
            (i-1)x_1 + x_2 &= 0 \\
            2x_1 + (-1 - i)x_2 &= 0
          \end{cases}
          \Rightarrow x_2 = (1 - i)x_1
          \Rightarrow X_1 = t\begin{pmatrix}
            1 \\ 1 - i
          \end{pmatrix}
        \]
        For $\lambda = -1$:
        \[
          A + I = \begin{pmatrix}
            i + 1 & 1 \\
            2 & 1 - i
          \end{pmatrix}
          \Rightarrow \begin{cases}
            (i+1)x_1 + x_2 &= 0 \\
            2x_1 + (1 - i)x_2 &= 0
          \end{cases}
          \Rightarrow x_2 = -(i + 1)x_1 
          \Rightarrow X_2 = t\begin{pmatrix}
            1 \\ -1 - i
          \end{pmatrix}
        \]
        \item Since $X_1$ and $X_2$ are linearly independent, 
        \[
          \braces{\begin{pmatrix}
            1 \\ 1 - i
          \end{pmatrix}, \begin{pmatrix}
            1 \\ -1 - i
          \end{pmatrix}}
        \]
        is a basis for $\CC^2$.

        \item It follows that 
        \[
          Q = \begin{pmatrix}
            1 & 1 \\
            1 - i & -1 - i
          \end{pmatrix}
          \mand 
          D = \begin{pmatrix}
            1 & 0 \\
            0 & -1
          \end{pmatrix}
        \]
        so $Q^{-1}$ exists since $\det(Q) = -2 \neq 0$ and $Q^{-1}AQ = D$.
      \end{enumerate}

      \part \begin{enumerate}[label=\roman*.)]
        \item The characteristic polynomial is 
        \begin{align*}
          \det(A - \lambda I) = \begin{vmatrix}
            2 - \lambda & 0 & -1 \\
            4 & 1 - \lambda & -4 \\
            2 & 0 & -\lambda - 1
          \end{vmatrix}
          &= (1 - \lambda)\Big((2 - \lambda)(-\lambda - 1) + 2\Big) = 0 \\
          &\Rightarrow\, (1 - \lambda)(\lambda^2 - \lambda) = 0 \\
          &\Rightarrow\, -\lambda(\lambda-1)^2 = 0 \\
          &\Rightarrow\, \lambda = 0, 1 \mul{2}
        \end{align*}
        Hence $A$ has eigenvalues $\lambda = 0, 1$. 

        \item For $\lambda = 0$:
        \[
          A - 0I = \begin{pmatrix}
            2 & 0 & -1 \\
            4 & 1 & -4 \\
            2 & 0 & -1
          \end{pmatrix}
          \xrightarrow[r_2 \to r_2 - 2r_1]{r_3 \to r_3 - r_1}
          \begin{pmatrix}
            2 & 0 & -1 \\
            0 & 1 & -2 \\
            0 & 0 & 0
          \end{pmatrix}
          \Rightarrow \begin{cases}
            x_3 &= t \\
            x_2 - 2x_3 &= 0 \\
            2x_1 - x_3 &= 0 
          \end{cases}
          \Rightarrow X_1 = t\begin{pmatrix}
            1/2 \\ 2 \\ 1
          \end{pmatrix}
        \]
        For $\lambda = 1$:
        \[
          A - I = \begin{pmatrix}
            1 & 0 & -1 \\
            4 & 0 & -4 \\
            2 & 0 & -2 
          \end{pmatrix}
          \to 
          \begin{pmatrix}
            1 & 0 & -1 \\
            0 & 0 & 0 \\
            0 & 0 & 0
          \end{pmatrix}
          \Rightarrow \begin{cases}
            x_3 &= t \\
            x_2 &= s \\
            x_1 - x_3 &= 0
          \end{cases}
          \Rightarrow X_2 = t\begin{pmatrix}
            1 \\ 0 \\ 1
          \end{pmatrix}
          + s\begin{pmatrix}
            0 \\ 1 \\ 0
          \end{pmatrix}
        \]

        \item Since eigenvectors corresponding to distinct eigenvalues are linearly independent (and the two vectors spanning the eigenspace for $\lambda = 1$ are independent), 
        \[
          \braces{\begin{pmatrix}
            1/2 \\ 2 \\ 1
          \end{pmatrix}, \begin{pmatrix}
            1 \\ 0 \\ 1
          \end{pmatrix}, \begin{pmatrix}
            0 \\ 1 \\ 0
          \end{pmatrix}}
        \]
        is a basis for $\RR^3$.

        \item Hence 
        \[
          Q = \begin{pmatrix}
            1/2 & 1 & 0 \\
            2 & 0 & 1 \\
            1 & 1 & 0
          \end{pmatrix}
          \mand 
          D = \begin{pmatrix}
            0 & 0 & 0 \\
            0 & 1 & 0 \\
            0 & 0 & 1
          \end{pmatrix}
        \]
        and $\det(Q) = 0.5 \neq 0$ so $Q^{-1}$ exists and $D = Q^{-1}AQ$.
      \end{enumerate}
    \end{parts}
  \end{problem}

  \begin{problem}{5.1.5f}
    Let $\gamma$ be the standard ordered basis for $P_3(\RR)$ so 
    \[
      [T]_\gamma = \begin{pmatrix}
        1 & 0 & 0 & 0 \\
        1 & 3 & 4 & 8 \\
        0 & 0 & 1 & 0 \\
        0 & 0 & 0 & 1
      \end{pmatrix}.
    \]
    Then the characteristic polynomial is 
    \[
      \det([T]_\gamma - \lambda I) = \begin{vmatrix}
        1 - \lambda & 0 & 0 & 0 \\
        1 & 3 - \lambda & 4 & 8 \\
        0 & 0 & 1 - \lambda & 0 \\
        0 & 0 & 0 & 1 - \lambda 
      \end{vmatrix}
      = (1 - \lambda)\begin{vmatrix}
        3 - \lambda & 4 & 8 \\
        0 & 1 - \lambda & 0 \\
        0 & 0 & 1 - \lambda 
      \end{vmatrix}
      = (1 - \lambda)^3(3 - \lambda) = 0
    \]
    so $T$ has eigenvalues $1 \mul{3}$ and $3$. For $\lambda = 1$:
    \[
      [T]_\gamma - I = \begin{pmatrix}
        0 & 0 & 0 & 0 \\
        1 & 2 & 4 & 8 \\
        0 & 0 & 0 & 0 \\
        0 & 0 & 0 & 0 
      \end{pmatrix}
      \Rightarrow \begin{cases}
        x_4 &= t \\
        x_3 &= r \\
        x_1 + 2x_2 + 4x_3 + 8x_4 &= 0 \Rightarrow x_2 = -4t - 2r - \frac{1}{2}s\\
        x_1 &= s
      \end{cases}
    \]
    Hence 
    \[
      X = s\begin{pmatrix}
        1 \\ -\frac{1}{2} \\ 0 \\ 0 
      \end{pmatrix}
      + r\begin{pmatrix}
        0 \\ -2 \\ 1 \\ 0
      \end{pmatrix}
      + t\begin{pmatrix}
        0 \\ -4 \\ 0 \\ 1
      \end{pmatrix}.
    \]
    Now, for $\lambda = 3$: 
    \[
      [T]_\gamma - 3I = \begin{pmatrix}
        -2 & 0 & 0 & 0 \\
        1 & 0 & 4 & 8 \\
        0 & 0 & -2 & 0 \\
        0 & 0 & 0 & -2 
      \end{pmatrix}
      \to \begin{pmatrix}
        1 & 0 & 0 & 0 \\
        1 & 0 & 4 & 8 \\
        0 & 0 & 1 & 0 \\
        0 & 0 & 0 & 1
      \end{pmatrix}
      \to \begin{pmatrix}
        1 & 0 & 0 & 0 \\
        0 & 0 & 0 & 0 \\
        0 & 0 & 1 & 0 \\
        0 & 0 & 0 & 1
      \end{pmatrix}
      \Rightarrow \begin{cases}
        x_4 &= 0 \\
        x_3 &= 0 \\
        x_2 &= t \\
        x_1 &= 0
      \end{cases}
    \]
    Hence, 
    \[
      X = t\begin{pmatrix}
        0 \\ 1 \\ 0 \\ 0
      \end{pmatrix}
    \]
    Thus, the ordered basis is $\beta = \braces{1 - \frac{x}{2}, x^2 - 2x, x^3 - 4x, x}$ and 
    \[
      [T]_\beta = \begin{pmatrix}
        1 & 0 & 0 & 0 \\
        0 & 1 & 0 & 0 \\
        0 & 0 & 1 & 0 \\
        0 & 0 & 0 & 3
      \end{pmatrix}
    \]
    is clearly diagonal.
  \end{problem}

  \begin{problem}{5.1.9}
    \begin{parts}
      \part \begin{proof}
        $(\Rightarrow)$ By way of the contrapositive assume that $0$ is an eigenvalue of $T$. It follows by definition that $\det(A - 0I) = \det(A) = 0$ where $A = [T]_\beta$ for any ordered basis $\beta$. However, it then follows that $A$ is not invertible by the Corollary to Thm~4.7 and so $T$ cannot be invertible by Thm~2.18. \\
        $(\Leftarrow)$ Conversely, suppose that $0$ is not an eigenvalue of $T$ and, for the sake of contradiction, that $T$ is not invertible. Then $A = [T]_\beta$ is not invertible for any ordered basis $\beta$ by Thm~2.18, so $\det(A) = 0$ by the Corollary to Thm~4.7. Hence $\det(A - 0I) = 0$, so $0$ is a root of the characteristic polynomial of $T$ and thus an eigenvalue of $T$, a contradiction. Therefore $T$ is invertible.
      \end{proof}

      \part \begin{proof}
        $(\Rightarrow)$ Let $T: V \to V$ be an invertible linear operator such that $\lambda^{-1} \in \FF$ is an eigenvalue of $T^{-1}$ (here $\lambda \neq 0$ by part $(a)$). By definition it follows that there exists some nonzero $w \in V$ such that 
        \[
          T^{-1}(w) = \lambda^{-1}w \Rightarrow T(T^{-1}(w)) = T(\lambda^{-1}w) \Rightarrow w = \lambda^{-1}T(w) \Rightarrow \lambda w = T(w)
        \]
        and hence $\lambda$ must be an eigenvalue of $T$. \\
        $(\Leftarrow)$ Conversely we may follow an almost identical argument assuming $\lambda$ is an eigenvalue for $T$ with eigenvector $w \neq 0$ (note $\lambda \neq 0$ by part $(a)$): 
        \[
          T(w) = \lambda w \Rightarrow T^{-1}(T(w)) = T^{-1}(\lambda w) \Rightarrow w = \lambda T^{-1}(w) \Rightarrow \lambda^{-1} w = T^{-1}(w)
        \]
        and hence $\lambda^{-1}$ is an eigenvalue for $T^{-1}$.
      \end{proof}
    \end{parts}
  \end{problem}

  \begin{problem}{5.1.10}
    \begin{proof}
      Let $M \in \MM_{n \times n}(\FF)$ be any upper triangular matrix. Then, 
      \[
        \det(M - \lambda I) = \begin{vmatrix}
          a_{11} - \lambda & a_{12} & a_{13} & \cdots & a_{1n} \\
          0 & a_{22} - \lambda & a_{23} & \cdots & a_{2n}  \\
          0 & 0 & a_{33} - \lambda & \cdots & a_{3n} \\
          \vdots & \vdots & \vdots & \ddots & \vdots \\
          0 & 0 & 0 & \cdots & a_{nn} - \lambda
        \end{vmatrix}
        = (a_{11} - \lambda)(a_{22} - \lambda) \cdots (a_{nn} - \lambda)
      \]
      using the fact that the determinant of any upper triangular matrix is the product of its diagonal entries (Exercise~$4.2.23$). Hence, the eigenvalues of $M$ are exactly its diagonal entries.
    \end{proof}
  \end{problem}

  \newpage 

  \begin{problem}{5.2.2}
    \begin{parts}
      \setcounter{partno}{3}
      \part We start by finding the characteristic polynomial 
      \begin{align*}
        \det(A - \lambda I) = \begin{vmatrix}
          7 - \lambda & -4 & 0 \\
          8 & -5 - \lambda & 0 \\
          6 & -6 & 3 - \lambda
        \end{vmatrix}
        &= (3 - \lambda)\Big((7 - \lambda)(-5 - \lambda) + 32\Big) \\
        &= (3 - \lambda)(\lambda^2 - 2\lambda - 3) \\
        &= (3 - \lambda)(\lambda + 1)(\lambda - 3) = 0 \\
        &\Rightarrow\, \lambda = 3 \mul{2}, -1
      \end{align*}
      So clearly the characteristic polynomial splits over $\RR$. For $\lambda_1 = 3$: 
      \[
        A - 3I = \begin{pmatrix}
          4 & -4 & 0 \\
          8 & -8 & 0 \\
          6 & -6 & 0 
        \end{pmatrix}
        \to \begin{pmatrix}
          1 & -1 & 0 \\
          0 & 0 & 0 \\
          0 & 0 & 0 
        \end{pmatrix}
        \Rightarrow \begin{cases}
          x_3 &= t \\
          x_2 &= r \\
          x_1 &= x_2
        \end{cases}
        \Rightarrow X = r\begin{pmatrix}
          1 \\ 1 \\ 0
        \end{pmatrix}
        + t\begin{pmatrix}
          0 \\ 0 \\ 1
        \end{pmatrix}
      \]
      thus 
      \[
        E_{\lambda_{1}} = \braces{
          r\begin{pmatrix}
            1 \\ 1 \\ 0
          \end{pmatrix}
          + t\begin{pmatrix}
            0 \\ 0 \\ 1
          \end{pmatrix}
          \mid r, t \in \RR
        }.
      \]
      For $\lambda_2 = -1$: 
      \[
          A + I = \begin{pmatrix}
            8 & -4 & 0 \\
            8 & -4 & 0 \\
            6 & -6 & 4
          \end{pmatrix}
          \to \begin{pmatrix}
            8 & -4 & 0 \\
            0 & 0 & 0 \\
            6 & -6 & 4
          \end{pmatrix}
          \Rightarrow \begin{cases}
            6x_1 - 6x_2 + 4x_3 &= 0 \\
            x_2 &= 2t \\
            8x_1 - 4x_2 &= 0
          \end{cases}
          \Rightarrow X = t\begin{pmatrix}
            1 \\ 2 \\ 3/2
          \end{pmatrix}
      \]
      thus 
      \[
        E_{\lambda_{2}} = \braces{
          t\begin{pmatrix}
            1 \\ 2 \\ 3/2
          \end{pmatrix}
          \mid t \in \RR
        }.
      \]
      Since $\dim(E_{\lambda_{1}}) = 2 = $ multiplicity of $\lambda_1 = 3$ and $\dim(E_{\lambda_{2}}) = 1 = $ multiplicity of $\lambda_2 = -1$, $A$ is diagonalizable. So, 
      \[
          Q = \begin{pmatrix}
            1 & 0 & 1 \\
            1 & 0 & 2 \\
            0 & 1 & 3/2
          \end{pmatrix}
          \mand 
          D = \begin{pmatrix}
            3 & 0 & 0 \\
            0 & 3 & 0 \\
            0 & 0 & -1
          \end{pmatrix}
      \] 

      \setcounter{partno}{5}
      \part Since $A$ is upper triangular we know that its eigenvalues are its diagonal entries: $1 \mul{2}$ and $3$. Since $1, 3 \in \RR$ the characteristic polynomial must split over $\RR$. For $\lambda_1 = 1$: 
      \[
          A - I = \begin{pmatrix}
            0 & 1 & 0 \\
            0 & 0 & 2 \\
            0 & 0 & 2
          \end{pmatrix}
          \to \begin{pmatrix}
            0 & 1 & 0 \\
            0 & 0 & 1 \\
            0 & 0 & 0
          \end{pmatrix}
          \Rightarrow \begin{cases}
            x_3 &= 0 \\
            x_2 &= 0 \\
            x_1 &= t 
          \end{cases}
          \Rightarrow X = t\begin{pmatrix}
            1 \\ 0 \\ 0
          \end{pmatrix}
      \]
      so 
      \[
          E_{\lambda_{1}} = \braces{
            t\begin{pmatrix}
              1 \\ 0 \\ 0
            \end{pmatrix}
            \mid t \in \RR
          }.
      \]
      Since $\dim(E_{\lambda_{1}}) = 1 \neq 2 = $ multiplicity of $\lambda_1 = 1$, $A$ is not diagonalizable.
    \end{parts}
  \end{problem}

  \begin{problem}{5.2.3d}
    Let $\alpha$ be the standard ordered basis of $P_2(\RR)$. Then 
    \begin{align*}
      [T]_\alpha = \begin{pmatrix}
        1 & 0 & 0 \\
        1 & 1 & 1 \\
        1 & 1 & 1
      \end{pmatrix} = A
      \Rightarrow \det(A - \lambda I) = \begin{vmatrix}
        1 - \lambda & 0 & 0 \\
        1 & 1 - \lambda & 1 \\
        1 & 1 & 1 - \lambda
      \end{vmatrix}
      &= (1 - \lambda)\Big((1 - \lambda)^2 - 1\Big) \\
      &= (1 - \lambda)(1 - 2\lambda + \lambda^2 - 1) \\
      &= (1 - \lambda)(\lambda^2 - 2\lambda) = 0 \\
      &\Rightarrow \lambda = 0, 1, 2
    \end{align*}
    Since $0, 1, 2 \in \RR$ the characteristic polynomial splits over $\RR$. Now, for $\lambda_1 = 0$: 
    \[
      A \to \begin{pmatrix}
        1 & 0 & 0 \\
        0 & 1 & 1 \\
        0 & 0 & 0
      \end{pmatrix}
      \Rightarrow \begin{cases}
        x_3 &= t \\
        x_2 &= -t \\
        x_1 &= 0 
      \end{cases}
      \Rightarrow E_{\lambda_{1}} = \braces{
        t\begin{pmatrix}
          0 \\ -1 \\ 1
        \end{pmatrix}
        \mid t \in \RR
      }.
    \]
    For $\lambda_2 = 1$:
    \[
        A - I = \begin{pmatrix}
          0 & 0 & 0 \\
          1 & 0 & 1 \\
          1 & 1 & 0 
        \end{pmatrix}
        \to \begin{pmatrix}
          1 & 1 & 0 \\
          1 & 0 & 1 \\
          0 & 0 & 0
        \end{pmatrix}
        \to \begin{pmatrix}
          1 & 1 & 0 \\
          0 & -1 & 1 \\
          0 & 0 & 0 
        \end{pmatrix}
        \Rightarrow \begin{cases}
          x_3 &= t \\
          x_2 &= t \\
          x_1 &= -t 
        \end{cases}
        \Rightarrow E_{\lambda_{2}} = \braces{
          t\begin{pmatrix}
            -1 \\ 1 \\ 1
          \end{pmatrix}
          \mid t \in \RR
        }.
    \]
    For $\lambda_3 = 2$: 
    \begin{align*}
      A - 2I &= \begin{pmatrix}
        -1 & 0 & 0 \\
        1 & -1 & 1 \\
        1 & 1 & -1 
      \end{pmatrix}
      \to \begin{pmatrix}
        1 & 0 & 0 \\
        1 & -1 & 1 \\
        1 & 1 & -1 
      \end{pmatrix}
      \to \begin{pmatrix}
        1 & 0 & 0 \\
        0 & -1 & 1 \\
        0 & 1 & -1 
      \end{pmatrix}
      \to \begin{pmatrix}
        1 & 0 & 0 \\
        0 & -1 & 1 \\
        0 & 0 & 0
      \end{pmatrix} \\
      &\Rightarrow\, \begin{cases}
        x_3 &= t \\
        x_2 &= t \\
        x_1 &= 0 
      \end{cases}
      \Rightarrow E_{\lambda_{3}} = \braces{
        t\begin{pmatrix}
          0 \\ 1 \\ 1
        \end{pmatrix}
        \mid t \in \RR
      }.
    \end{align*}
    Hence since the dimension of each eigenspace equals the multiplicity of its respective eigenvalue, $T$ is diagonalizable. Thus, for 
    \[
      \beta = \braces{x^2 - x, x^2 + x - 1, x + x^2}
    \]
    $[T]_\beta = \operatorname{diag}(0, 1, 2)$ is diagonal.
  \end{problem}

  \begin{problem}{5.2.13}
    \begin{parts}
      \part \begin{proof}
        Consider an invertible linear operator $T$  and let $\lambda$ be one eigenvalue of $T$ so $\lambda^{-1}$ is an eigenvalue of $T^{-1}$. Consider the eigenspace of $T$ corresponding to $\lambda$, $E_{\lambda}$, and the eigenspace of $T^{-1}$ corresponding to $\lambda^{-1}$, $E'_{\lambda^{-1}}$ (note that $\lambda \neq 0$ by Exercise~5.1.9). Now, let $v \in E_\lambda$ be any vector in the eigenspace of $\lambda$, so that $T(v) = \lambda v$. Then, 
        \[
          T(v) = \lambda v \Rightarrow T^{-1}(T(v)) = T^{-1}(\lambda v) \Rightarrow v = \lambda T^{-1}(v) \Rightarrow \lambda^{-1} v = T^{-1}(v)
        \]
        so $v \in E'_{\lambda^{-1}}$ and since $v$ was arbitrary, $E_{\lambda} \subseteq E'_{\lambda^{-1}}$. An almost identical argument can be made by flipping $\lambda$ and $\lambda^{-1}$, $T$ and $T^{-1}$, and $E_\lambda$ and $E'_{\lambda^{-1}}$ to show $E'_{\lambda^{-1}} \subseteq E_\lambda$. Hence, $E_\lambda = E'_{\lambda^{-1}}$ as was to be shown. 
      \end{proof}

      \part \begin{proof}
        It follows from $(a)$ that every eigenvector of $T$ is an eigenvector of $T^{-1}$, with eigenvalue replaced by its reciprocal. Suppose that $T$ is diagonalizable. Then there is a basis $\beta$ of $V$ consisting of eigenvectors of $T$, and by $(a)$ each vector of $\beta$ is also an eigenvector of $T^{-1}$. Hence $\beta$ is a basis of $V$ consisting of eigenvectors of $T^{-1}$, so $T^{-1}$ is diagonalizable.
      \end{proof}
    \end{parts}
  \end{problem}

  \begin{problem}{5.2.21}
    \begin{proof}
      $(\Rightarrow)$ Suppose that $V$ is the direct sum of $W_1, \dots, W_k$ and let $\beta_i$ be a basis for $W_i$. We claim that the list $\beta$ obtained by concatenating $\beta_1, \dots, \beta_k$ is a basis for $V$. It spans $V$ since $\sum W_i = V$. For independence, suppose a linear combination of the vectors in $\beta$ equals $0$ and let $w_i \in W_i$ be the part of that combination coming from $\beta_i$, so $w_1 + \cdots + w_k = 0$. Then for each $j$,
      \[
        w_j = -\sum_{i \neq j} w_i \in W_j \cap \sum_{i \neq j} W_i = \braces{0},
      \]
      so every $w_j = 0$, and since each $\beta_i$ is linearly independent, every coefficient is $0$. Hence $\beta$ is a basis for $V$. In particular $\beta$ has no repeated vectors, so
      \[
        \dim(V) = |\beta| = \sum_{i=1}^{k} |\beta_i| = \sum_{i=1}^{k} \dim(W_i).
      \] \\
      $(\Leftarrow)$ Now suppose that $\dim(V) = \sum_{i=1}^{k} \dim(W_i)$ and let $\beta$ be the list obtained by concatenating $\beta_1, \dots, \beta_k$ where each $\beta_i$ is a basis for $W_i$. Since $\sum W_i = V$, clearly $\beta$ spans $V$, and since $\beta$ has exactly $\sum_{i=1}^{k} \dim(W_i) = \dim(V)$ vectors, $\beta$ must be a basis for $V$. In particular $\beta$ is linearly independent. Now, consider $x \in W_j \cap \sum_{i \neq j} W_i$ so
      \[
        x = a_1 v_1 + \cdots + a_m v_m \mand x = b_1 u_1 + \cdots + b_p u_p
        \Rightarrow 0 = (a_1 v_1 + \cdots + a_m v_m) - (b_1 u_1 + \cdots + b_p u_p)
      \]
      for $v_1, \dots, v_m \in \beta_j$ and $u_1, \dots, u_p \in \bigcup_{i \neq j} \beta_i$ and scalars $a_1, \dots, a_m, b_1, \dots, b_p \in \FF$. The vectors $v_1, \dots, v_m, u_1, \dots, u_p$ are distinct vectors of $\beta$, so since $\beta$ is linearly independent all of the scalars must be $0$. In particular $a_1 = \cdots = a_m = 0$ and so $x = 0$. Thus $W_j \cap \sum_{i \neq j} W_i = \braces{0}$ for every $j$, and so $V$ is the direct sum of $W_1, \dots, W_k$, as was to be shown.
    \end{proof}
  \end{problem}

  \begin{problem}{5.2.22}
    \begin{proof}
      Define $W_i = \Span{\beta_i}$ so $W_i \subseteq V$. Since $\beta_i \cap \beta_j = \varnothing$ for $i \neq j$, then $|\beta| = \sum|\beta_i|$ but $|\beta| = \dim(V)$, and $|\beta_i| = \dim(W_i)$ since $\beta_i$ is linearly independent (being a subset of the basis $\beta$) and spans $W_i$, so 
      \[
        \dim(V) = |\beta| = \sum|\beta_i| = \sum\dim(W_i).
      \]
      As well, clearly 
      \[
        \sum_{i = 1}^{k} W_i = V
      \]
      since $\beta_i$ partition $\beta$ so by exercise $(5.2.21)$, 
      \[
        V = W_1 \oplus \cdots \oplus W_k = \Span{\beta_1} \oplus \cdots \oplus \Span{\beta_k}
      \]
      as was to be shown.
    \end{proof}
  \end{problem}
\end{document}